\documentclass[11pt,letterpaper]{article}
\usepackage[T1]{fontenc}
\usepackage{amsmath,amsthm}
\usepackage{newtxtext,newtxmath}
\usepackage[margin=1.04in,top=0.88in,bottom=0.90in]{geometry}
\usepackage{microtype}
\usepackage[hidelinks,pdfdisplaydoctitle]{hyperref}
\hypersetup{pdftitle={Solving a dense linear system has the same arithmetic exponent as matrix multiplication},
  pdfsubject={Exact solution of a rational linear system with one right-hand side}}
\newtheorem{theorem}{Theorem}
\newtheorem{lemma}[theorem]{Lemma}
\newcommand{\Q}{\mathbb{Q}}
\newcommand{\rankq}{\underline{R}_{\Q}}
\newcommand{\floor}[1]{\lfloor #1\rfloor}
\setlength{\parindent}{1.15em}
\setlength{\emergencystretch}{1em}
\raggedbottom
\frenchspacing

\begin{document}
\begin{center}
  {\LARGE\bfseries Solving a dense linear system has the same\\
    arithmetic exponent as matrix multiplication\par}
\end{center}
\vspace{0.3em}

\begin{abstract}
Solving a nonsingular rational \(n\times n\) system \(Ax=b\) with one right-hand
side has the arithmetic exponent \(\omega\) of square matrix multiplication.
We prove that the same exponent is necessary even when \(b=e_1\) is fixed
in advance and the only output is \(x_1\). From any exact straight-line
program of \(s\) gates for this scalar we construct a division-free circuit
that computes a nonzero multiple of the generic determinant with at most
\(3s\) multiplications. Andrews's bound for such multiples gives
\(s\ge\floor{n/4}^{\omega}/18\) for every \(n\ge8\). In particular, the
value \(\omega=9/4\) reported in the OpenAI preprint applies even to a
single solution coordinate.
\end{abstract}

\section{One right-hand side}

Given a nonsingular \(n\times n\) matrix \(A\) over \(\Q\) and one vector
\(b\in\Q^n\), compute \(x=A^{-1}b\). Computing the inverse would solve the system
for all right-hand sides. Can one do better with just one? Allaire and
Kaber list the equivalence with matrix multiplication for the full map
\((A,b)\mapsto A^{-1}b\) \cite[Theorem~4.4.1]{allaire}. We ask how much
of the output already forces that exponent.

Let \(S(n)\) be the minimum gate count for the full solution, and \(C(n)\)
the minimum for \(A\mapsto(A^{-1})_{11}\), the first coordinate when
\(b=e_1\). Each gate performs one of \(+,-,\times,/\) on two available
scalars. Rational constants and reuse are allowed. The gate sequence is
fixed, and all divisions must be defined on every admissible input. Write
\(\sigma\) and \(\sigma_1\) for the respective exponents, where the exponent
of a cost \(H\) is \(\inf\{t:H(n)=O(n^t)\}\). Let \(\omega\) denote the
arithmetic exponent of square matrix multiplication over \(\Q\).

\begin{theorem}\label{thm:main}
The full solution and the single prescribed coordinate satisfy
\[
                       \sigma=\sigma_1=\omega .
\]
In each dimension \(n\ge8\), \(C(n)\ge\floor{n/4}^{\omega}/18\).
\end{theorem}

\begin{proof}
Fixing \(b=e_1\) in a full solver and retaining its first output gives
\(C(n)\le S(n)\). Put \(k=\floor{n/4}\). A program for the coordinate also
works in dimension \(4k\) by feeding it \(\operatorname{diag}(Y,I_{n-4k})\).
Lemma~\ref{lem:denominator} then gives a nonzero multiple of \(\det Y\)
using at most \(3C(n)\) multiplications. Write \(\rankq(k)\) for the border
rank of \(k\times k\) matrix multiplication over \(\Q\), with approximations
over \(\Q(\varepsilon)\) as in \cite[\S2.5]{andrews}. By Andrews's
theorem on the determinant ideal \cite[Theorem~3.4]{andrews},
\[
              \frac{\rankq(k)}{6}\le 3C(n).
\]
Bini's theorem identifies \(\omega\) with the growth exponent of
\(\rankq(k)\) \cite{bini} (see also \cite[Lemma~2.15]{andrews}).
For fixed \(k\ge2\), tensoring a border decomposition
and padding to the next power of \(k\) gives
\(\rankq(m)\le\rankq(k)^{\lceil\log_k m\rceil}\). Thus
\(\omega\le\log_k\rankq(k)\). This proves the finite bound and
\(\sigma_1\ge\omega\).

For the upper bound, fix \(\tau>\omega\). Form \(G=A^\top A\) and return
\(G^{-1}(A^\top b)\). The matrix \(G\) is positive definite, as are its
leading principal block and the associated Schur complement.
After padding \(G\) with an identity block to a power of two, recursive
block inversion \cite{bunch} therefore uses a fixed sequence of legal
divisions. The inversion cost satisfies \(T(n)\le2T(n/2)+O(n^\tau)\),
using fast matrix multiplication for the block products. Thus
\(S(n)=O(n^\tau)\) and \(\sigma\le\omega\).
\end{proof}

\newpage
\section{The denominator of one coordinate}

Let \(X=(x_{ij})\) be a matrix of independent variables, \(D=\det X\), and
\(\Delta\) the cofactor of \(x_{11}\). Cramer's rule gives
\((X^{-1})_{11}=\Delta/D\). The reduced denominator of this one entry is
already the full determinant, up to a nonzero constant. The next lemma
bounds the cost of recovering a multiple of it from an arbitrary program.

\begin{lemma}\label{lem:denominator}
Let \(n\ge2\). If a straight-line program of \(s\) gates returns
\((A^{-1})_{11}\) on every invertible rational \(n\times n\) input, then a
division-free circuit with at most \(3s\) multiplication gates computes a
nonzero element of the ideal \((D)\subset\Q[x_{ij}]\).
\end{lemma}

\begin{proof}
Represent each register by an unreduced pair of polynomials, initially
\((x_{ij},1)\) or \((c,1)\). For operands represented by \((p,q)\)
and \((r,t)\), the new register is
\[
\begin{array}{c@{\qquad}c@{\qquad}c}
  +\text{ or }- & \times & /\\[2pt]
  (pt\pm rq,\,qt) & (pr,\,qt) & (pt,\,qr).
\end{array}
\]
Each gate costs at most three multiplications. Both components of every pair
remain available for reuse. Let \((P,Q)\) be the final pair. On invertible
rational inputs every intermediate denominator is nonzero by induction, so
\(Q(I)\ne0\).

Correctness gives \(PD-Q\Delta=0\) on invertible rational matrices.
Multiplying by \(D\) gives an identity on \emph{all} rational matrices,
including singular ones. Since \(\Q\) is infinite and \(D\ne0\), it follows
that \(PD=Q\Delta\) in \(\Q[x_{ij}]\).

Suppose \(g\) divides both \(D\) and \(\Delta\). Degrees in a given variable
add under multiplication of nonzero polynomials. Since \(\Delta\) omits the
first row, so does \(g\). Replace that row by \(e_j^\top\). This fixes \(g\)
and turns \(D\) into the nonzero signed minor omitting column \(j\). Its
divisor \(g\) is therefore independent of column \(j\). This holds for
every \(j\), so \(g\) is constant. Unique factorization and \(PD=Q\Delta\)
now yield \(D\mid Q\).
\end{proof}

\paragraph{The value \(9/4\).}
Substituting the value of \(\omega\) in the OpenAI preprint
\cite[Theorem~1.1]{openai} gives \(C(n)\ge \floor{n/4}^{9/4}/18
\ge n^{9/4}/2304\) for \(n\ge8\), since \(\floor{n/4}\ge n/8\).
Permuting rows and columns gives the same bound for any prescribed entry
of \(A^{-1}\). It also applies to the exact optimal value of
\(\min\{x_1:Ax=e_1,\ x\in\Q^n\}\), whose feasible set is a singleton.

\paragraph{Formal verification.}
The Lean~4.19.0 proof verifies the denominator construction starting from
correctness on all nonsingular rational inputs. Its numerical lower bound
takes Andrews's determinant-ideal bound and the pointwise border-rank
bound in \cite[Corollary~1.2]{openai} as hypotheses. At revision
\texttt{b81e394}, \texttt{./scripts/verify} builds the proof and audits its axioms.

\begin{thebibliography}{9}
\setlength{\itemsep}{0.12em}
\bibitem{allaire} G.~Allaire and S.~M. Kaber.
\href{https://doi.org/10.1007/978-0-387-68918-0}{\emph{Numerical Linear Algebra}}.
Springer, 2008, Theorem~4.4.1.
\bibitem{andrews} R.~Andrews. \emph{On matrix multiplication and polynomial
identity testing}. FOCS 2022. \href{https://arxiv.org/abs/2208.01078}{arXiv:2208.01078}.
\bibitem{bini} D.~Bini. \href{https://doi.org/10.1007/BF02575865}{\emph{Relations between
exact and approximate bilinear algorithms. Applications}}. Calcolo 17 (1980), 87--97.
\bibitem{bunch} J.~R. Bunch and J.~E. Hopcroft. \href{https://doi.org/10.1090/S0025-5718-1974-0331751-8}{\emph{Triangular factorization and inversion by fast matrix multiplication}}.
Mathematics of Computation 28 (1974), 231--236.
\bibitem{openai} OpenAI. \emph{The matrix multiplication exponent is nine fourths}.
\href{https://github.com/openai/math/blob/main/preprints/The-matrix-multiplication-exponent-is-nine-fourths-November-24-2026/paper.pdf}{Preprint}, 24 November 2026.
\end{thebibliography}
\end{document}
